% GENERATED FILE. Edit canonical lessons, then run npm run generate:books.
% canonical-source-hash: fnv1a64:5bd8a67a7c836d7b
% canonical-lessons: AR-C380-adala, AR-C380-musawah, AR-C380-zalim, AR-C380-adil, AR-C380-istahaqqa

\chapter{Justice and fairness}
\label{ch:ar-justice-and-fairness}
\hlchaptermodality{380}

\begin{chapteropening}
\textbf{By the end of this chapter:} \emph{I can talk about justice and equality, say what is fair or unjust, and say what someone deserves.}

Complete the last lesson of chapter 380: I can talk about justice and equality, say what is fair or unjust, and say what someone deserves.
\end{chapteropening}

\section[\texorpdfstring{\ar{عدالة} (ʿadāla)}{عدالة (ʿadāla)}]{\ar{عدالة} (ʿadāla) — justice}

\label{lesson:AR-C380-adala}



\begin{quote}
Before the new one: say the Arabic for an activist, then the Arabic for democracy.
\end{quote}

\subsection*{You'll want to know: \ar{عدالة}}
\textbf{\ar{عدالة}} — \emph{ʿadāla} — \textquotedblleft{}justice\textquotedblright{}.

\begin{quote}
\ar{يطالب} \ar{الناس} \ar{بالعدالة}.
\end{quote}

\begin{quote}
\emph{yuṭālibu an-nās bi-l-ʿadāla.}
\end{quote}

\begin{quote}
People are demanding justice.
\end{quote}

\begin{cousinweb}
\textbf{\ar{عدالة}} (\emph{ʿadāla}) is built on the native Arabic root \emph{ʿ-d-l}, \textquotedblleft{}to be equal, straight\textquotedblright{} (Lane). \emph{ʿAdāla} is a verbal noun, being fair and upright: justice.
\end{cousinweb}

\begin{grammarlens}[title={What you've built}]
One more word for giving reasons and opinions.
\end{grammarlens}

\subsection*{Your turn}
\begin{itemize}
\raggedright
  \item \emph{Say it:} \emph{ʿadāla}
  \item \emph{Say it:} \emph{ʿadāla}, once more
  \item \emph{Say it:} \emph{yuṭālibu an-nās bi-l-ʿadāla.} — People are demanding justice.
  \item \emph{Recall:} say the Arabic for an activist, then the Arabic for democracy, then say \emph{ʿadāla} again
\end{itemize}

\begin{tcolorbox}[breakable,colback=teal!4,colframe=teal!35!black,title={Before you move on}]
What does \ar{عدالة} mean? (\textquotedblleft{}Justice\textquotedblright{}.) Say it once more.
\end{tcolorbox}
\section[\texorpdfstring{\ar{مساواة} (musāwāh)}{مساواة (musāwāh)}]{\ar{مساواة} (musāwāh) — equality}

\label{lesson:AR-C380-musawah}



\begin{quote}
Before the new one: say the Arabic for democracy, then the Arabic for justice.
\end{quote}

\subsection*{You'll want to know: \ar{مساواة}}
\textbf{\ar{مساواة}} — \emph{musāwāh} — \textquotedblleft{}equality\textquotedblright{}.

\begin{quote}
\ar{نريد} \ar{المساواة} \ar{بين} \ar{الرجل} \ar{والمرأة}.
\end{quote}

\begin{quote}
\emph{nurīdu al-musāwāh bayna ar-rajul wa-l-marʾa.}
\end{quote}

\begin{quote}
We want equality between men and women.
\end{quote}

\begin{cousinweb}
\textbf{\ar{مساواة}} (\emph{musāwāh}) is built on the native Arabic root \emph{s-w-y}, \textquotedblleft{}to be equal, level\textquotedblright{} (Lane). \emph{Musāwāh} is the Form III verbal noun, making equal: equality.
\end{cousinweb}

\begin{grammarlens}[title={What you've built}]
One more word for giving reasons and opinions.
\end{grammarlens}

\subsection*{Your turn}
\begin{itemize}
\raggedright
  \item \emph{Say it:} \emph{musāwāh}
  \item \emph{Say it:} \emph{musāwāh}, once more
  \item \emph{Say it:} \emph{nurīdu al-musāwāh bayna ar-rajul wa-l-marʾa.} — We want equality between men and women.
  \item \emph{Recall:} say the Arabic for democracy, then the Arabic for justice, then say \emph{musāwāh} again
\end{itemize}

\begin{tcolorbox}[breakable,colback=teal!4,colframe=teal!35!black,title={Before you move on}]
What does \ar{مساواة} mean? (\textquotedblleft{}Equality\textquotedblright{}.) Say it once more.
\end{tcolorbox}
\section[\texorpdfstring{\ar{ظالم} (ẓālim)}{ظالم (ẓālim)}]{\ar{ظالم} (ẓālim) — unjust}

\label{lesson:AR-C380-zalim}



\begin{quote}
Before the new one: say the Arabic for justice, then the Arabic for equality.
\end{quote}

\subsection*{You'll want to know: \ar{ظالم}}
\textbf{\ar{ظالم}} — \emph{ẓālim} — \textquotedblleft{}unjust\textquotedblright{}.

\begin{quote}
\ar{كان} \ar{القرار} \ar{ظالما} \ar{للعمال}.
\end{quote}

\begin{quote}
\emph{kāna al-qarār ẓāliman li-l-ʿummāl.}
\end{quote}

\begin{quote}
The decision was unjust to the workers.
\end{quote}

\begin{cousinweb}
\textbf{\ar{ظالم}} (\emph{ẓālim}) is built on the native Arabic root \emph{ẓ-l-m}, \textquotedblleft{}to put a thing in the wrong place\textquotedblright{} (Lane). \emph{Ẓālim} is the active participle: one who does wrong, unjust.
\end{cousinweb}

\begin{grammarlens}[title={What you've built}]
One more word for giving reasons and opinions.
\end{grammarlens}

\subsection*{Your turn}
\begin{itemize}
\raggedright
  \item \emph{Say it:} \emph{ẓālim}
  \item \emph{Say it:} \emph{ẓālim}, once more
  \item \emph{Say it:} \emph{kāna al-qarār ẓāliman li-l-ʿummāl.} — The decision was unjust to the workers.
  \item \emph{Recall:} say the Arabic for justice, then the Arabic for equality, then say \emph{ẓālim} again
\end{itemize}

\begin{tcolorbox}[breakable,colback=teal!4,colframe=teal!35!black,title={Before you move on}]
What does \ar{ظالم} mean? (\textquotedblleft{}Unjust\textquotedblright{}.) Say it once more.
\end{tcolorbox}
\section[\texorpdfstring{\ar{عادل} (ʿādil)}{عادل (ʿādil)}]{\ar{عادل} (ʿādil) — fair}

\label{lesson:AR-C380-adil}



\begin{quote}
Before the new one: say the Arabic for equality, then the Arabic for unjust.
\end{quote}

\subsection*{You'll want to know: \ar{عادل}}
\textbf{\ar{عادل}} — \emph{ʿādil} — \textquotedblleft{}fair\textquotedblright{}.

\begin{quote}
\ar{كان} \ar{المعلم} \ar{عادلا} \ar{مع} \ar{كل} \ar{الطلاب}.
\end{quote}

\begin{quote}
\emph{kāna al-muʿallim ʿādilan maʿa kulli aṭ-ṭullāb.}
\end{quote}

\begin{quote}
The teacher was fair to all the students.
\end{quote}

\begin{cousinweb}
\textbf{\ar{عادل}} (\emph{ʿādil}) is built on the native Arabic root \emph{ʿ-d-l}, \textquotedblleft{}to be equal, straight\textquotedblright{} (Lane). \emph{ʿĀdil} is the active participle: one who acts fairly, fair-minded.
\end{cousinweb}

\begin{grammarlens}[title={What you've built}]
One more word for giving reasons and opinions.
\end{grammarlens}

\subsection*{Your turn}
\begin{itemize}
\raggedright
  \item \emph{Say it:} \emph{ʿādil}
  \item \emph{Say it:} \emph{ʿādil}, once more
  \item \emph{Say it:} \emph{kāna al-muʿallim ʿādilan maʿa kulli aṭ-ṭullāb.} — The teacher was fair to all the students.
  \item \emph{Recall:} say the Arabic for equality, then the Arabic for unjust, then say \emph{ʿādil} again
\end{itemize}

\begin{tcolorbox}[breakable,colback=teal!4,colframe=teal!35!black,title={Before you move on}]
What does \ar{عادل} mean? (\textquotedblleft{}Fair\textquotedblright{}.) Say it once more.
\end{tcolorbox}
\section[\texorpdfstring{\ar{استحق} (istaḥaqqa)}{استحق (istaḥaqqa)}]{\ar{استحق} (istaḥaqqa) — to deserve}

\label{lesson:AR-C380-istahaqqa}



\begin{quote}
Before the new one: say the Arabic for unjust, then the Arabic for fair.
\end{quote}

\subsection*{You'll want to know: \ar{استحق}}
\textbf{\ar{استحق}} — \emph{istaḥaqqa} — \textquotedblleft{}to deserve\textquotedblright{}.

\begin{quote}
\ar{يستحق} \ar{هذا} \ar{العامل} \ar{راتبا} \ar{أعلى}.
\end{quote}

\begin{quote}
\emph{yastaḥiqqu hādhā al-ʿāmil rātiban aʿlā.}
\end{quote}

\begin{quote}
This worker deserves a higher salary.
\end{quote}

Present: \ar{يستحق}, \emph{yastaḥiqq}.

\begin{cousinweb}
\textbf{\ar{استحق}} (\emph{istaḥaqqa}) is built on the native Arabic root \emph{ḥ-q-q}, \textquotedblleft{}to be true, to be due\textquotedblright{} (Lane). It is Form X, to seek what is due: to deserve.
\end{cousinweb}

\begin{grammarlens}[title={What you've built}]
One more word for giving reasons and opinions.
\end{grammarlens}

\subsection*{Your turn}
\begin{itemize}
\raggedright
  \item \emph{Say it:} \emph{istaḥaqqa}
  \item \emph{Say it:} \emph{istaḥaqqa}, once more
  \item \emph{Say it:} \emph{yastaḥiqqu hādhā al-ʿāmil rātiban aʿlā.} — This worker deserves a higher salary.
  \item \emph{Recall:} say the Arabic for unjust, then the Arabic for fair, then say \emph{istaḥaqqa} again
\end{itemize}

\begin{tcolorbox}[breakable,colback=teal!4,colframe=teal!35!black,title={Before you move on}]
What does \ar{استحق} mean? (\textquotedblleft{}To deserve\textquotedblright{}.) Say it once more.
\end{tcolorbox}
